\section{Planting the wiring at marked points}\label{subsec:planting}

A station in the input may have wires to many other stations. To move it
as one insertion, we encode its ends of those wires in a single leg.
The sphere form must then test whether the bits at opposite ends agree.
We implement that test separately in each occupied cell of each network.

There are two sphere forms in a cell. The form $A$ carries the station
legs and the seams to neighboring cells. Extra marks on $A$ supply a
joint state that tests the wiring. The private form $B$ realizes this
joint state through a separate insertion table at each extra mark.
Thus, after planting, each insertion has just one leg on each sphere,
along with any exposed or control bits. At original stations and seams
the $B$ leg is the vacuum. Figure~\ref{fig:stack} shows the two kinds
of insertions.

\begin{figure}[htb]\centering
\begin{tikzpicture}[x=1cm,y=1cm,font=\small,
  strip/.style={draw=ink,rounded corners=3pt,line width=.65pt},
  ins/.style={draw=ink,rounded corners=2pt,fill=white,inner xsep=7pt,inner ysep=4pt},
  leg/.style={draw=ink,line width=.8pt}]
\draw[strip,fill=ink!10] (-4.35,.98) rectangle (4.35,1.55);
\draw[strip,fill=ink!5] (-4.35,-1.55) rectangle (4.35,-.98);
\node[anchor=west] at (-4.22,1.26) {Sewn sphere \(A\)};
\node[anchor=east] at (4.22,1.26) {\(R\) species};
\node[anchor=west] at (-4.22,-1.26) {Private sphere \(B\)};
\node[anchor=east] at (4.22,-1.26) {\(R/b\) species};
\draw[leg] (-5.35,1.26)--(-4.35,1.26);
\draw[leg] (4.35,1.26)--(5.35,1.26);
\node[above,font=\footnotesize] at (-4.85,1.26) {seam};
\node[above,font=\footnotesize] at (4.85,1.26) {seam};
\node[ins] (sta) at (-2.15,0) {original station};
\node[ins] (aux) at (2.0,0) {auxiliary insertion};
\draw[leg] (sta.north)--(-2.15,.98);
\draw[leg,densely dotted] (sta.south)--(-2.15,-.98);
\node[anchor=west,font=\footnotesize] at (-2.1,-.6) {vacuum};
\draw[leg] (aux.north)--(2.0,.98);
\draw[leg] (aux.south)--(2.0,-.98);
\node[anchor=west] at (2.05,.62) {\(z^P\)};
\node[anchor=west] at (2.05,-.62) {\(z\)};
\foreach \x/\y in {-2.15/.98,2.0/.98,2.0/-.98}
    \fill[ink] (\x,\y) circle (1.5pt);
\draw[draw=ink,fill=white] (-2.15,-.98) circle (1.5pt);
\end{tikzpicture}
\caption{The station table encodes its incident internal wire bits on \(A\).
The auxiliary tables couple \(A\) to \(B\) and together impose equality
along the wires. Only \(A\) is sewn to neighboring cells.}\label{fig:stack}
\end{figure}

A mark at $z$ on $B$ corresponds to the mark at $z^P$ on $A$, where $P$
is a power of the characteristic. Moving $z$ by $t$ moves $z^P$ by $t^P$.
In units of the Laurent order on $A$, the energy of the pair is therefore
$e_A+P^{-1}e_B$. A given energy budget allows $P$ times as much energy
on $B$. We keep the corresponding leg narrow by using fewer species on
$B$. This will bound the width left behind when a station crosses a seam.

Fix $g,\ell$ as in Lemma~\ref{lem:local}. Let $F$ be the first
characteristic power at least $D^{1000(\ell+1)}$ and let $a=\lceil F^3\rceil$.
For a chunk count $M$ and a sufficiently large characteristic power $b>F$,
take $R$ to be the largest multiple of $bM$ at most $T/a$. Set
\begin{equation}\label{eq:plant-scales}
 R'=R/b,\qquad P=b/F,\qquad \kappa=P^{-1},\qquad H=Ra^2 .
\end{equation}
Here $R$ and $R'$ count the species on $A$ and $B$, respectively, and
$T/(2a)\le R\le T/a$ for sufficiently large $T$. If $e_A+\kappa e_B\le h$,
Equation~\eqref{eq:local-width} bounds the widths by constant multiples
of $\sqrt{Rh}$ and $\sqrt{R'Ph}=\sqrt{Rh/F}$. In bits, up to rounding,
\begin{equation}\label{eq:plant-widths}
 w_A=O\bigl(T\sqrt{h/H}\bigr),\qquad
 w_B=O\bigl(T\sqrt{h/(FH)}\bigr).
\end{equation}

\subsection{The wiring form for one species}

Combine the wires between each pair of distinct stations into one ordered
bundle, absorbing endpoint permutations into the station tables. Contract
self-loops in those tables first. Split each internal bundle into packets
of $R$ bits, padding the last packet with zeros. Species $s$ will carry
bit $s$ of every packet. A station meets $O_g(D)$ bundles, and a cell
together with its incident edges meets $O_g(D^2)$. At a station of
arity $t_i$, one species therefore carries at most $t_i/R+O_g(D)$ bits.

Choose distinct wedge basis states at an integer energy $h_i$ to encode
the bit strings at station $i$. The logarithm of the single-species
dimension is asymptotic to $(\pi/\sqrt3)\sqrt{h_i}$ through integer
levels by Lemma~\ref{lem:dimension}. Since $R\ge T/(2a)$ and $a\ge D$,
we may take
\[
                 h_i=O_g\bigl(a^2(1+t_i/T)^2\bigr).
\]
At seam $j$ the packets across the interface contribute
$O_g(DT/R+D^2)=O_g(Da)$ bits to one species. An integer energy
$q_j=O_g(D^2a^2)$ suffices. The two cells use the same bit order and
the same basis codes there. For fixed $D$ the packet counts per
species, energies and codes range over finite sets.

Define the multilinear form $L$ on these station and seam levels by
its coefficients. Give a basis tuple coefficient one if all codes are
valid and both ends of every bit agree. For a wire to another cell,
compare its station end with its seam end. Every other coefficient
is zero. We will sew a multispecies seam at a fixed \emph{total}
energy. To recover the correct energy in each individual species,
put $Q=\sum_j q_j$ and extend $L$ by zero to all other tuples of
nonnegative half-integer seam energies of sum at most $Q$. The station
energies stay fixed.

\begin{lemma}[Auxiliary marks of bounded weight]\label{lem:plant-probes}
Inserting a joint state at polynomially many extra marks on a
single-species sphere realizes the extended form $L$. At each extra
mark the state uses only the vacuum and $\psi_0$. In every word in its
support the number of occupied marks is even and at most $O_g(Da)$.
\end{lemma}
\begin{proof}
A wedge of length $d$ has energy at least
$\tfrac12+\tfrac32+\cdots+(d-\tfrac12)=d^2/2$. At the stations the
total length is $O_g(a\sum_i(1+t_i/T))=O_g(Da)$. At the at most $g$
seams the same bound follows from Cauchy--Schwarz and their total
energy bound $Q=O_g(D^2a^2)$. List all pairs $(v,m)$ of original
marks and modes occurring in the resulting input spaces. The list
has polynomial size in $D,a$.

An insertion of $\psi_0$ at a new point $z$ pairs with those modes by
the covector with coordinates $(z-z_v)^{-m-1}$. Distinct principal
parts give linearly independent rational functions in every
characteristic. Choose distinct regular points whose covectors span
the dual mode space. Indeed, a nonzero linear combination annihilated
by the points already chosen is nonzero at some unused point of the
algebraic closure. At most one point per listed mode is needed, all
in a finite extension. Add vacua if necessary to obtain at least
three auxiliary marks.

For a subset of $d$ new marks, insert $\psi_0$ there and the vacuum
at the others. Consider first the part of the resulting Pfaffian
in which every auxiliary generator pairs with an original generator.
This part is itself a linear combination of insertion tests on at
most $d$ marks. Indeed, order all auxiliary generators last and group
the matchings by their exact set of internal auxiliary pairs. Each
nonempty set contributes its contracted scalar, with a sign depending
only on those paired positions, times the part on the remaining marks
with no internal auxiliary pairs. Subtracting these contributions
inductively isolates the desired part.

Now take a nonzero vector in the even part of the direct sum of input
spaces, and let $d$ be its greatest nonzero total wedge length. A test
on $d$ marks with no internal auxiliary pairs vanishes on every shorter
wedge. On the length-$d$ component it evaluates the exterior product of
the chosen $d$ covectors, up to a common sign. The original basis
tuples are distinct basis wedges in the exterior algebra on the list
of pairs $(v,m)$. Some choice of $d$ covectors therefore detects the
component. The resulting even tests span the full dual of the even
part. Each actual insertion used has even weight at most $d=O_g(Da)$
and vanishes on odd vectors. The prescribed integer energies give
even lengths, so the zero extension of $L$ also vanishes on the odd
part. A linear combination of insertion tests now gives the required
joint state.
\end{proof}

Use one copy of this state in each of the $R$ species. At station $i$
the insertion table enforces level $h_i$ in every species, decodes the
packets, checks that padded bits are zero and supplies the original
table entry, with its exposed bits. At seam $j$ we sew by basis
equality at total energy $Rq_j$. To see why this recovers the original
wiring, group the Pfaffian by species. Order the original and auxiliary
marks in two consecutive lists and reorder each list by species.
The auxiliary reordering sign can be absorbed in the joint state.
Interleaving the lists costs no further sign because every species
has even auxiliary weight. The original reordering sign depends only on
the input wedges. In any nonzero product after the auxiliary state has
been contracted, the seam tuple of each species
is either the prescribed tuple or has sum greater than $Q$, by the zero
extension. The sums over all species add to $RQ$, so every species
has the prescribed tuple. All the original wedges then have even
length and the remaining reordering sign is one. The product tests
exactly the equality of the bits along each wire. We also place one
vacuum collector in each cell for the motion step.

\subsection{Realizing the joint auxiliary state}\label{subsec:plant-private}

Lemma~\ref{lem:plant-probes} prescribes one tensor shared by all
auxiliary marks. For motion we need a separate insertion table at
each mark. The private sphere supplies the shared tensor through
local maps. One $B$ species will serve a block of $b$ species on $A$.

Partition the $R$ sewn species into $R'$ consecutive blocks. A block
has \emph{type} $\tau=(d_i)_{i=1}^m$ if $d_i$ of its modes $\psi_0$
are occupied at auxiliary mark $i$. Projecting the block tensor onto
these local subspaces gives a tensor $J_\tau$, including the sign
for reordering within the block. Its $i$th leg has dimension
$\binom b{d_i}$. For a fixed ordered word of block types, the joint
tensor is the product of the $J_\tau$ up to a known sign. Exchanging
two site lists from different blocks contributes only the product of
their lengths modulo two to the sign exponent. Retain the types with
$J_\tau\ne0$.

Write $H_2$ for binary entropy in bits, with $H_2(0)=H_2(1)=0$.
Round each $H_2(d_i/b)$ up to a strictly larger point $u_i$ of a
fixed positive grid of mesh $1/m$. Then
\begin{equation}\label{eq:plant-ent}
 \log_2\binom b{d_i}\le b u_i,\qquad
 \max_i u_i\le 2,\qquad \sum_i u_i=O_{g,\ell}(Da\log D).
\end{equation}
Indeed, Lemma~\ref{lem:plant-probes} gives $s=\sum_i d_i/b=O_g(Da)$.
Concavity bounds the sum of the entropies by $s\log_2(em/s)$,
with value zero at $s=0$. The $m$ marks are polynomially many in $D$
for fixed $\ell$ and characteristic, and rounding adds at most one
to the sum.

At the corresponding $B$ marks, take the direct sum $S_\tau$ of all
nonzero single-species sphere forms with energy at most $C(bu_i)^2$
at site $i$, where $C$ is an absolute constant chosen below. For
sufficiently large $b$ there is a constant $c_\tau$ such that for
every $d\ge1$ a product of local linear maps gives
\begin{equation}\label{eq:plant-comparison}
 \bigoplus_{s=1}^{c_\tau} S_\tau^{\otimes d}
                    \longrightarrow J_\tau^{\otimes d}.
\end{equation}
The independence of $c_\tau$ from $d$ will let one selector index
serve every occurrence of the type.

For the comparison, take natural logarithms and work over the
algebraic closure, where the exponent is still strictly below $9/4$.
For a normalized spectral price at $x_1,\ldots,x_m$ put
$M_u=\inf_o\sum_i u_i\,d(o,x_i)$. Preparing $J_\tau$ at $o$ and
transporting its legs bounds its price by $b(\log 2)\sum_i u_i d(o,x_i)$
by Equation~\eqref{eq:plant-ent}. Taking the infimum gives $b(\log 2)M_u$. Applying the same transport bound to
the size-two diagonal, whose price is $\log 2$, gives
$\inf_o\sum_i d(o,x_i)\ge1$. Hence $M_u\ge\min_i u_i>0$.
Apply Corollary~\ref{cor:universal} at scale $N=b$ with width caps
$c_0\sqrt C\,u_i$. Its comparison tensor is precisely $S_\tau$,
whose price is $b(c_0\sqrt C\,M_u+o(1))$ uniformly in the normalized
price.
Choose $c_0\sqrt C>2\log 2$. For fixed $D$ there are finitely many
grid vectors and choices of $A$ positions, all defined over a fixed
finite extension. Taking $P$th roots permutes that field, so the
possible $B$ positions also range over a finite set as $b$ varies.
It follows that $S_\tau$ strictly dominates $J_\tau$ in every price
for all sufficiently large characteristic powers $b$.
The field-extension and strict comparison results of
Section~\ref{sec:prices-background} give Equation~\eqref{eq:plant-comparison}
over the fixed field of definition. The maps extend to the field of
the input tables.

Fix the type word. For each type of multiplicity $d>0$, expand the
restriction in Equation~\eqref{eq:plant-comparison} by choosing one
copy index $s$ and one $B$-energy tuple for each of the $d$ factors.
Assign the factors to private species in the original block order.
Their product is the multispecies sphere form up to the sign for
permuting the homogeneous lists. Each site map now supplies a local
table coupling the $A$ and $B$ legs at that mark. It may act jointly
on all blocks of a given type \emph{at that mark}. Summing the choices
and type words produces the required auxiliary tensor. The $B$ legs
at the original stations and seams stay in the vacuum.

\paragraph{Control and storage.}
A \emph{private selector} records the ordered type word, the
per-block $B$ energies and one index $s$ per occurring type. A single
table checks validity, including zero padding of unused slots, and
supplies the permutation sign. The types and twice the energies cost
$O(m\log(2b))$ bits per block, since $d_i\le b$ and the doubled
energies are $O(b^2)$. The copy indices cost $O_{D,b}(1)$ bits in
total. With $R/b$ blocks, choosing first $b$ and then $T$ large makes
the selector shorter than $\epsilon T$ for any fixed $\epsilon>0$.

On $A$ every occupied auxiliary mode costs energy $1/2$, giving
total auxiliary energy $O_g(RDa)$. On one $B$ species the sum of
the caps is $O(b^2\sum_i u_i^2)=O_{g,\ell}(b^2 Da\log D)$ by
Equation~\eqref{eq:plant-ent}. Multiplying by $R'$ and converting to
$A$ units gives the cell bound
\[
 \sum_i(e_{A,i}+\kappa e_{B,i})
    =O_{g,\ell}(RDa+F RDa\log D)\le D H
\]
for sufficiently large $D$, since $H=Ra^2$ and $a\ge F^3$.
For storage, use the actual spaces of the selected block. At mark $i$
either leg needs only $O(bu_i+1)$ bits, by
Equations~\eqref{eq:plant-ent} and~\eqref{eq:local-width}.
Over all blocks, each auxiliary insertion needs $O(R)=O(T/a)$ bits
for these two legs. Summed over the marks, this costs
$O_{g,\ell}(DT\log D)$, taking $b\ge m$ to absorb rounding.

The selector is summed inside the network, but the outer layout must
know the energy and width at each mark. Split the blocks into $M$ equal
consecutive chunks. In the outer sum, specify for each chunk the
histogram of its pairs (type, $B$-energy tuple). A histogram records
the multiplicities of the pairs, while the private selector records
their order. The alphabet is fixed independently of $T$, so the
histograms cost $O_{D,b,M}(\log(T+1))$ bits per cell. They fix the
energy and the concatenated code length of each leg in each chunk.
The table checks the histograms against the selector and uses its block
order to decode. Each valid selector occurs in exactly one outer term.
Thus $M$ can be chosen for the chunk layout, then $\epsilon$ for its
selector routing, and finally $b$ for the comparison and selector width,
all before $T$ grows.
